\section{De la structure complexe interne à la géométrie projective de Bloch}
\label{sec:geometria-proyectiva-bloch-rev2}
\index{sphère de Bloch}
\index{structure complexe!réalification orthogonale}
\index{groupe icosaédrique binaire}
\index{Möbius!action projective}

La racine interne de moins l'identité possède deux réalisations coordonnées. Dans
la Partie I, l'incidence de Paley--Witt construit
\(J_W^2=-I\) sur \(\mathbb F_3^6\), organise le code ternaire et réalise
la structure \(\mathbb F_9^3\) ; dans une réalisation analytique, un opérateur
orthogonal \(J_E\) avec \(J_E^2=-I\) transforme un espace réel de dimension
paire en un espace complexe. Cette seconde flèche ajoute topologie, norme et
complétion : elle conserve la relation quadratique, mais n'identifie pas les corps
de scalaires.

Le passage complet est typé par
\begin{equation}
 \boxed{
 \begin{gathered}
 C_P\longrightarrow J_W^2=-I_6
 \longrightarrow \mathbb F_3^6\simeq\mathbb F_9^3,\\
 (E_4,J_E),\quad \dim_{\mathbb R}E_4=4,
 \longrightarrow
 \mathbb P(\mathcal H_J)\simeq\mathbb {CP}^{1}
 \longrightarrow S^2,\\
 \mathrm{SU}(2)\longrightarrow\mathrm{PSU}(2)\simeq\mathrm{SO}(3)
 \subset\mathrm{PSL}(2,\mathbb C)\cong\text{\rmfamily Möb}^{+}(S^2),\\
 2A_5=\pi^{-1}(A_5)\subset\mathrm{SU}(2).
 \end{gathered}}
 \label{eq:cadena-paley-bloch-rev2}
\end{equation}
La première ligne conserve ses preuves dans les
\cref{chap:paley-codigo-witt,thm:f9-tres-desde-paley} ; on construit ici les
deux lignes projectives.

\subsection{Réalification de la phase et espace des rayons}

Soit \((E,g)\) un espace euclidien réel de dimension paire et soit
\(J_E\in\operatorname{End}(E)\) tel que
\[
 J_E^2=-I,
 \qquad g(J_Eu,J_Ev)=g(u,v).
\]
Alors \(J_E^{\mathsf T}=-J_E\) et
\[
 R_J(\theta):=e^{\theta J_E}
 =\cos\theta\,I+\sin\theta\,J_E
\]
est orthogonal, avec
\(R_J(\theta)R_J(\phi)=R_J(\theta+\phi)\). L'action
\((a+ib)u:=au+bJ_Eu\) réalise la carte complexe sans utiliser le scalaire
\(i\) comme antécédent de l'opérateur.

Soit maintenant \(E_4\subseteq E\) un sous-espace \(J_E\)-stable de dimension réelle
quatre et soit \(\mathcal H_J=(E_4,J_E)\) l'espace de Hilbert résultant, de
dimension complexe deux.
Un vecteur unitaire
\[
 \psi=\begin{pmatrix}z_0\\z_1\end{pmatrix},
 \qquad |z_0|^2+|z_1|^2=1,
\]
détermine un rayon \([\psi]\in\mathbb {CP}^{1}\). Définissons
\begin{equation}
 r_1=2\operatorname{Re}(\overline z_0z_1),
 \qquad
 r_2=2\operatorname{Im}(\overline z_0z_1),
 \qquad
 r_3=|z_0|^2-|z_1|^2.
 \label{eq:coordenadas-bloch-rev2}
\end{equation}

\begin{theorem}[Identification projective de Bloch]
\label{thm:bloch-proyectivo-rev2}
L'application
\[
 \mathcal B:\mathbb {CP}^{1}\longrightarrow S^2,
 \qquad [\psi]\longmapsto(r_1,r_2,r_3),
\]
est une bijection lisse. Le projecteur de rang un associé au rayon est
\begin{equation}
 \rho_\psi=|\psi\rangle\langle\psi|
 =\frac12\bigl(I_2+r_1\sigma_1+r_2\sigma_2+r_3\sigma_3\bigr),
 \label{eq:proyector-bloch-rev2}
\end{equation}
où \(\sigma_1,\sigma_2,\sigma_3\) sont les matrices de Pauli.
\end{theorem}

\begin{proof}
La phase globale laisse invariantes les trois coordonnées et
\[
 r_1^2+r_2^2+r_3^2
 =4|z_0|^2|z_1|^2+(|z_0|^2-|z_1|^2)^2=1.
\]
Réciproquement, chaque \(\boldsymbol r\in S^2\) détermine le projecteur positif
\(\rho=(I+\boldsymbol r\cdot\boldsymbol\sigma)/2\), qui satisfait
\(\rho^2=\rho\) et \(\operatorname{tr}\rho=1\) ; son image est un unique rayon.
Les cartes affines de \(\mathbb {CP}^{1}\) rendent lisses l'application et son inverse
\cite{Bloch1946}.
\end{proof}

La suite de Hopf
\begin{equation}
 S^1\longrightarrow S^3\longrightarrow S^2
 \label{eq:fibrado-hopf-bloch-rev2}
\end{equation}
expose la perte publiée par la projectivisation : le vecteur normalisé
appartient à \(S^3\), l'observable projectif publie un point de \(S^2\) et
la fibre \(S^1\) conserve la phase commune que la carte identifie
\cite{Hopf1931}. C'est une réalisation précise de la différence HMT entre
état enrichi et coordonnée visible.

\subsection{Rotations, transformations de Möbius et relèvement spinoriel}

Tout élément de \(\mathrm{SU}(2)\) peut s'écrire
\[
 U=\begin{pmatrix}
 \alpha&\beta\\-\overline\beta&\overline\alpha
 \end{pmatrix},
 \qquad |\alpha|^2+|\beta|^2=1.
\]

\begin{theorem}[Action spinorielle sur la sphère de Bloch]
\label{thm:su2-so3-bloch-rev2}
Pour chaque \(U\in\mathrm{SU}(2)\), il existe une unique rotation
\(R_U\in\mathrm{SO}(3)\) telle que
\[
 U(\boldsymbol r\cdot\boldsymbol\sigma)U^\dagger
 =(R_U\boldsymbol r)\cdot\boldsymbol\sigma.
\]
L'application \(U\mapsto R_U\) est un homomorphisme surjectif de noyau
\(\{\pm I_2\}\).
\end{theorem}

\begin{proof}
La conjugaison conserve l'espace réel tridimensionnel des matrices
hermitiennes de trace nulle et le produit scalaire
\(\langle A,B\rangle=\tfrac12\operatorname{tr}(AB)\), si bien qu'elle induit une
rotation. Si elle agit trivialement, \(U\) commute avec les trois matrices de Pauli
et est scalaire ; le déterminant un laisse \(U=\pm I_2\). La représentation
axe--angle fournit les deux images réciproques de chaque rotation.
\end{proof}

Dans la carte affine \(z=z_0/z_1\) de \(\mathbb {CP}^{1}\), la même action est
la transformation de Möbius
\[
 z\longmapsto
 \frac{\alpha z+\beta}{-\overline\beta z+\overline\alpha}.
\]
La rotation de Bloch et la transformation fractionnelle sont deux coordonnées
d'une unique action projective : la première agit sur les projecteurs hermitiens et
la seconde sur la sphère de Riemann.

Le groupe des rotations de l'icosaèdre est \(A_5\subset\mathrm{SO}(3)\). Son
image réciproque définit le groupe icosaédrique binaire
\[
 2A_5:=\pi^{-1}(A_5)\subset\mathrm{SU}(2),
\]
et produit la suite exacte
\begin{equation}
 1\longrightarrow\{\pm I_2\}
 \longrightarrow2A_5\longrightarrow A_5\longrightarrow1.
 \label{eq:extension-central-2a5-rev2}
\end{equation}
En particulier, \(|2A_5|=120\). Pour une rotation d'ordre impair \(n\), le
relèvement de demi-angle possède l'ordre \(2n\), tandis que son opposé possède
l'ordre \(n\). Les rotations icosaédriques d'ordres trois et cinq admettent donc
des paires de relèvements d'ordres \((6,3)\) et \((10,5)\),
respectivement. Cette spécialisation relie la
géométrie projective des états à deux niveaux à la branche de McKay et aux
symétries exceptionnelles, dont les modules et les réseaux demeurent dans leurs
constructions de référence de la Partie I.

\begin{statusbox}[title={Provenance et portée de la chaîne projective}]
L'identité de Paley--Witt, la structure
\(\mathbb F_3^6\simeq\mathbb F_9^3\), l'action
\(\mathrm{SU}(2)\to\mathrm{SO}(3)\), la réalisation de Möbius et l'extension
\(2A_5\) sont des résultats démontrés dans leurs domaines. La chaîne
\eqref{eq:cadena-paley-bloch-rev2} est une réalisation fonctorielle unifiée :
elle déclare explicitement la réalisation qui change de scalaires et ne transforme pas
la racine finie \(J_W\) en un opérateur analytique sans cette donnée supplémentaire.
\end{statusbox}
